\section{Capítulo~\ref{sec:dofaalt}. Otras teorías de \nt{DOPA}}

Cada uno de los cuadros ampliados se apoya en sus propias medidas directas de la dopamina (dispersión de los puntos de inversión, respuestas a lo desagradable y a lo nuevo, crecimiento lento, respuestas al cambio de sabor), pero las fórmulas que las explican son teorías, y entre dos de ellas los experimentos se contradicen por ahora.

{\footnotesize
\setlength{\tabcolsep}{3pt}\renewcommand{\arraystretch}{1.15}
\begin{longtable}{>{\raggedright}p{0.5cm}>{\raggedright}p{4.0cm}>{\raggedright}p{5.6cm}>{\raggedright}p{2.3cm}>{\raggedright\arraybackslash}p{1.7cm}}
\toprule
N.º & Proposición & Experimento y qué se midió & Fuente & Estado \\
\midrule\endfirsthead
\toprule
N.º & Proposición & Experimento y qué se midió & Fuente & Estado \\
\midrule\endhead
1 & La regla asimétrica~\eqref{eq:asymmetric} converge al punto~\eqref{eq:expectile}; con $\kappa=1/2$ es la media; para una gota con probabilidad $1/2$, $V=\kappa$ (fig.~\ref{fig:expectiles}) & El punto~\eqref{eq:expectile} es un expectil, solución de un problema de mínimos cuadrados con pesos distintos para las desviaciones positivas y negativas; para el ejemplo del capítulo, la deducción está en el propio capítulo. & \cite{NeweyPowell1987}; deducción en el capítulo & teoría \\
2 & Distintas neuronas \nt{DOPA} tienen distintos puntos de inversión, ordenados según la asimetría (proposición~\ref{prop:expectile}) & Ratones, neuronas \nt{DOPA} del área tegmental ventral marcadas optogenéticamente, recompensas de distinto tamaño: el punto en que el pico pasa a valle difiere entre neuronas; es más alto en las neuronas con respuestas más asimétricas; a partir de las respuestas del grupo se reconstruye la forma de la distribución de recompensas. Que esta sea la función principal de la dispersión es una hipótesis. & \cite{Dabney2020} & medido \\
3 & Parte de las neuronas \nt{DOPA} responde a lo desagradable con un pico & Monos: unas neuronas \nt{DOPA} se excitan con la señal previa a la recompensa y se inhiben con la señal previa a un soplo de aire en el ojo; otras se excitan con ambas; los grupos están en partes distintas del mesencéfalo. & \cite{Matsumoto2009, BrombergMartin2010} & medido \\
4 & El módulo del error o la novedad fija la tasa de aprendizaje & En los trabajos citados no hay experimento directo en que la señal de importancia de las neuronas \nt{DOPA} cambie la tasa de aprendizaje; la revisión propone el papel de una señal de “atención, es importante”. & \cite{BrombergMartin2010} & hipótesis \\
5 & Un cable aparte para el eje del peligro & Ratones: las neuronas \nt{DOPA} que van al extremo posterior del estriado responden a estímulos nuevos e intensos, no a la recompensa; su destrucción debilita la evitación de un objeto nuevo. & \cite{Menegas2018} & medido \\
6 & Tiempo óptimo de la acción $\tau_a^*=\sqrt{C/\bar R}$~\eqref{eq:vigor} & Mínimo de la suma $C/\tau_a+\bar R\tau_a$; deducción en el capítulo. En el modelo original, la rapidez de las acciones la fija el precio del tiempo, igual a la recompensa media. & \cite{Niv2007}; deducción en el capítulo & teoría \\
7 & El nivel lento de dopamina lleva $\bar R$ y fija la rapidez de las acciones (proposición~\ref{prop:multiplex}) & Ratas, microdiálisis y voltamperometría en el núcleo accumbens: el nivel de dopamina sigue minuto a minuto la frecuencia de recompensas y la rapidez de las acciones. En humanos, la L-DOPA refuerza la dependencia del tiempo de reacción respecto de la frecuencia media de recompensas. Que el nivel lento sea precisamente $\bar R$ no está demostrado. & \cite{Hamid2016, Beierholm2013vigor} & hipótesis \\
8 & El nivel lento se forma a partir de dos fuentes: la liberación cambia en las salidas sin que cambie la frecuencia de espigas, pero también hay crecimiento lento en la propia frecuencia & Ratas, una misma tarea: la liberación en el núcleo accumbens sigue la expectativa cambiante de recompensa, y la frecuencia de espigas de las neuronas \nt{DOPA} identificadas del área tegmental ventral, no. Ratones en un pasillo virtual (fila~10): el crecimiento gradual al acercarse a la recompensa está también en las espigas de neuronas \nt{DOPA} identificadas. & \cite{Mohebi2019, Kim2020} & medido \\
9 & La dopamina crece de forma gradual mientras el animal se acerca a una recompensa lejana & Ratas en un laberinto, voltamperometría en el estriado: la concentración de dopamina sube en segundos a medida que se acercan a la meta; la pendiente depende de la distancia y del tamaño de la recompensa. & \cite{Howe2013ramp, Hamid2016} & medido \\
10 & El crecimiento es el error $\delta$ al afinarse la estimación~\eqref{eq:ramp}; un salto en el pasillo da un pico (fig.~\ref{fig:ramp}) & La fórmula y el número $\delta=0.75\,V$ por segundo son deducción del capítulo. Experimento: ratones en un pasillo virtual, traslado instantáneo más cerca de la meta; las espigas de los somas, el calcio de somas y salidas y la concentración en el estriado responden como un error, no como la expectativa $V$. Cuál de las dos explicaciones vale en todas las salidas se discute. & \cite{Kim2020} & medido \\
11 & Mirar atrás: la dopamina informa de una medida del vínculo causal~\eqref{eq:retro} & Ratones, liberación de dopamina en el núcleo accumbens en experimentos en que esta teoría y el error~\eqref{eq:ctd} predicen cosas distintas: en varios experimentos la liberación siguió a la primera. & \cite{Jeong2022} & hipótesis \\
12 & La respuesta de la dopamina durante el aprendizaje se desplaza poco a poco de la recompensa a la señal previa & Ratones, señal de las salidas de las neuronas \nt{DOPA} en el estriado ventral a lo largo del aprendizaje: el pico pasa gradualmente del momento de la recompensa al de la señal previa, como exige el aprendizaje según~\eqref{eq:ctd}. & \cite{Amo2022} & medido \\
13 & Las neuronas \nt{DOPA} responden al cambio de sabor de la recompensa con el mismo valor & Ratas, registro de neuronas del área tegmental ventral: sustituir un jugo por otro de igual valor provoca respuesta, aunque el error~\eqref{eq:ctd} es cero. & \cite{Takahashi2017} & medido \\
14 & Picos artificiales enseñan el vínculo entre dos estímulos neutros & Ratas, experimento de precondicionamiento sensorial con dos estímulos sin recompensa: la activación optogenética de las neuronas \nt{DOPA} en la transición del primero al segundo crea el vínculo, y su inhibición impide aprenderlo. & \cite{Sharpe2017} & medido \\
15 & Vector de errores por rasgos~\eqref{eq:vectortd} & Teoría del error de predicción de rasgos; no hay comprobación directa de la forma vectorial. & \cite{Gardner2018} & hipótesis \\
16 & Distintas neuronas \nt{DOPA} llevan en una misma tarea magnitudes distintas & Ratones en un pasillo virtual, imagen de dos fotones de neuronas \nt{DOPA}: unas se relacionan con la recompensa, otras con la velocidad y el giro, otras con las señales previas y la precisión de la elección. & \cite{Engelhard2019} & medido \\
17 & Un mazo de cables en lugar de un cable (proposición~\ref{prop:bundle}) & Resumen de las filas~2--16: cada descomposición está medida por separado; no hay un cuadro experimental único que muestre que todas son partes de una misma señal. & --- & hipótesis \\
\bottomrule
\end{longtable}}
